\begin{abstract}
Printed circuit board (PCB) interconnects introduce insertion loss, reflections
and crosstalk that become increasingly important at high data rates.
Electromagnetic simulations characterise these effects accurately but are
expensive for repeated design-space exploration. This work develops a staged
machine-learning surrogate workflow for one PCB topology, progressing from
regression baselines to curve-aware neural models and finally a
frequency-conditioned network that predicts the reciprocal complex 12-port
scattering matrix from ten design parameters over 200 frequency points. On
held-out designs, the final model achieved a complex mean absolute error of
0.0686, preserved reciprocity exactly, and showed no passivity violations on the
evaluated grid. The experiments also reveal a trade-off between complete
response modelling and selected-path accuracy, with sharp transmission nulls
remaining difficult to reproduce. The results demonstrate both the feasibility
and limitations of complete scattering-matrix surrogate modelling for
signal-integrity analysis.
\end{abstract}

\begin{IEEEkeywords}
scattering parameters, signal integrity, surrogate modelling, neural networks,
physical consistency
\end{IEEEkeywords}
